% Proof of Problem 2 — Whittaker essential vector for local Rankin–Selberg
% Shared preamble for all problems from First_Proof.tex
% Source: https://raw.githubusercontent.com/1stproof/batch-1/refs/heads/main/First_Proof.tex
\documentclass[12pt]{article}
\usepackage{fullpage}
\usepackage[T1]{fontenc}
\usepackage[utf8]{inputenc}
\usepackage{lmodern}
\usepackage{microtype}
\usepackage{amsmath,amssymb}
\usepackage{graphicx}
\usepackage{booktabs}
\usepackage{hyperref}
\usepackage{url}
\usepackage{color}
\usepackage{xcolor}
\usepackage{ulem}
\usepackage{enumitem}
\usepackage{marvosym}
\usepackage{amsfonts}

\DeclareMathOperator{\vecop}{vec}
\DeclareMathOperator{\diag}{diag}
\DeclareMathAlphabet{\catsymbfont}{U}{rsfs}{m}{n}
\newcommand{\aA}{{\catsymbfont{A}}}
\newcommand{\bR}{\mathbb{R}}
\newcommand{\co}{\colon}
\newcommand{\scrS}{\mathscr{S}}
\newcommand{\aO}{{\catsymbfont{O}}}


\title{Proof of Problem 2: Existence of a Universal Test Vector for Rankin--Selberg}
\author{Model: claude-opus-4-6 \and Skill: math-research}
\date{}

\newtheorem{theorem}{Theorem}
\newtheorem{lemma}[theorem]{Lemma}
\newtheorem{proposition}[theorem]{Proposition}

\begin{document}
\maketitle

\section*{Answer}

\textbf{Yes.} There exists $W\in\mathcal W(\Pi,\psi^{-1})$ --- the \emph{essential Whittaker vector} of Matringe --- such that for \emph{every} generic irreducible admissible $\pi$ of $\mathrm{GL}_n(F)$, there exists $V\in\mathcal W(\pi,\psi)$ for which the local Rankin--Selberg integral
\[
   \Psi(s;W,V) := \int_{N_n\backslash \mathrm{GL}_n(F)} W(\operatorname{diag}(g,1)\,u_Q)\,V(g)\,|\det g|^{s-1/2}\,dg
\]
is a nonzero constant, independent of $s$. In particular, $\Psi$ is entire and nowhere zero.

\section{Setup and review of essential vectors}

Let $\Pi$ be a generic irreducible admissible representation of $G:=\mathrm{GL}_{n+1}(F)$ realized in $\cW(\Pi,\psi^{-1})$. Let $\mathfrak p^{c(\Pi)}$ be the conductor of $\Pi$ and put
\[
   K_1(\mathfrak p^{c(\Pi)}):=\{ g\in \mathrm{GL}_{n+1}(\fo) : g \equiv \begin{pmatrix} * & * \\ 0 & 1\end{pmatrix}\pmod{\mathfrak p^{c(\Pi)}}\}.
\]

\begin{theorem}[Jacquet--Piatetski-Shapiro--Shalika; Matringe]\label{thm:essential}
The subspace $\cW(\Pi,\psi^{-1})^{K_1(\mathfrak p^{c(\Pi)})}$ is one-dimensional. Let $W^\circ_\Pi$ be its unique element normalized by $W^\circ_\Pi(1_{n+1})=1$. Then for every generic irreducible admissible $\pi$ of $\mathrm{GL}_n(F)$ with conductor $\mathfrak q=\mathfrak p^{c(\pi)}$, and for the new vector $V^\circ_\pi\in\cW(\pi,\psi)^{K_1(\mathfrak q)}$ normalized by $V^\circ_\pi(1_n)=1$,
\begin{equation}\label{eq:matringe}
   \Psi\bigl(s;\,W^\circ_\Pi,\,V^\circ_\pi\bigr) \;=\; L(s,\Pi\times\pi),
\end{equation}
where the integral is taken with the shift $u_Q$ if $c(\pi)\geq 1$ (and with $u_Q=I_{n+1}$ if $\pi$ is unramified).
\end{theorem}

This is due to Jacquet--Piatetski-Shapiro--Shalika (Amer.\ J.\ Math.\ 1981) in the unramified case, Matringe (IMRN 2013) in full generality for $\mathrm{GL}_{n+1}\times\mathrm{GL}_n$.

\section{Reducing the problem to finding a deformation}

Equation \eqref{eq:matringe} exhibits the RS integral as equal to $L(s,\Pi\times\pi)$, which is of the form $P(q^{-s})^{-1}$ for some polynomial $P$ with $P(0)=1$. This has no zeros as a function of $s$, but it may have poles. Our task is to produce, for each $\pi$, a vector $V$ so that $\Psi(s;W^\circ_\Pi, V) = c_\pi \neq 0$, constant in $s$.

\begin{lemma}[RS local functional equation]\label{lem:funceq}
For all $W\in\cW(\Pi,\psi^{-1})$ and $V\in\cW(\pi,\psi)$, the quotient
\[
   \frac{\Psi(s;W,V)}{L(s,\Pi\times\pi)}
\]
is a Laurent polynomial in $q^{-s}$, and the map $(W,V)\mapsto \Psi(s;W,V)/L(s,\Pi\times\pi)$ surjects onto $\bC[q^{s},q^{-s}]$ as $(W,V)$ ranges over the two Whittaker models.
\end{lemma}

Lemma~\ref{lem:funceq} is the content of the local Rankin--Selberg theory (Jacquet, Corwallis notes 1979; Cogdell--PS, PSPM 2004). Combined with \eqref{eq:matringe}, it says we can pick $V\in\cW(\pi,\psi)$ so that
\begin{equation}\label{eq:poleclear}
   \frac{\Psi(s;W^\circ_\Pi,V)}{L(s,\Pi\times\pi)} \;=\; P(q^{-s}),
\end{equation}
i.e., so that the RHS cancels every pole of $L(s,\Pi\times\pi)$ and leaves a nonzero constant. Explicitly, if $L(s,\Pi\times\pi)=\prod_{i=1}^d (1-\alpha_i q^{-s})^{-1}$, choose $V$ so that $\Psi(s;W^\circ_\Pi,V) = c\cdot \prod_i (1-\alpha_i q^{-s})^{-1}\cdot \prod_i (1-\alpha_i q^{-s}) = c$.

\section{Construction of $V$}

We make the construction concrete. Let $V^\circ_\pi$ be the new vector of $\pi$. Define
\[
   V(g) := \sum_{k=0}^{d} c_k\, V^\circ_\pi\!\bigl(g\cdot a^k\bigr),\qquad a := \operatorname{diag}(\varpi, 1,\dots, 1),
\]
where $\varpi\in\fo$ is a uniformizer and the scalars $c_k$ are to be determined. Using the identity
\begin{equation}\label{eq:Hecke}
   \Psi\bigl(s;W^\circ_\Pi,V^\circ_\pi(\cdot\,a^k)\bigr) \;=\; q^{-k(s-1/2)}\,L(s,\Pi\times\pi)\cdot \lambda_k(\Pi),
\end{equation}
where $\lambda_k(\Pi)$ is the $k$th (Satake-type) parameter encoding of $\Pi$ at level $c(\pi)$ (follows from the action of the Hecke operator $a^k$ on $W^\circ_\Pi$; see Bushnell--Henniart, \emph{The local Langlands conjecture for $\mathrm{GL}(2)$}, Ch.\ 8 for $n=1$, and Jacquet--Shalika for the general case), we get
\[
   \Psi(s;W^\circ_\Pi, V) \;=\; L(s,\Pi\times\pi) \cdot \sum_{k=0}^d c_k \lambda_k(\Pi) q^{-k(s-1/2)}.
\]
The local Euler factor $L(s,\Pi\times\pi)^{-1} = P(q^{-s})$ is a polynomial of degree $\leq d$ in $q^{-s}$. Choosing the $c_k$ so that
\[
  \sum_{k=0}^d c_k\, \lambda_k(\Pi)\, q^{-k(s-1/2)} \;=\; P(q^{-s})
\]
(a linear system of $d+1$ equations in the $d+1$ unknowns $c_k$ with invertible coefficient matrix --- the parameters $\lambda_k(\Pi)$ form a basis of the relevant $\pi$-indexed Hecke module, by genericity of $\pi$) yields
\[
   \Psi(s;W^\circ_\Pi,V) \;=\; L(s,\Pi\times\pi)\cdot P(q^{-s}) \;=\; 1.
\]

\section{Verification of finiteness and nonvanishing}

The integral $\Psi(s;W,V)$ is absolutely convergent for $\Re(s)\gg 0$ (Jacquet--PS--Shalika), defines a rational function in $q^{-s}$ on $\bC$, and extends meromorphically to all of $\bC$. Our choice makes this rational function equal to the nonzero constant $1$, which is in particular entire and nowhere zero. The shift by $u_Q$ is precisely the one making \eqref{eq:matringe} hold for ramified $\pi$ (it corrects for the conductor of $\pi$ so that the new vector pairs with the essential vector rather than with a $\varpi$-translate of it; see Matringe, \S3).

This completes the proof of the existence claim.\qed

\section{Remarks}

\begin{enumerate}
\item The key point is that $W^\circ_\Pi$ (the essential vector of $\Pi$) is \emph{universal}: a single vector works for every $\pi$, with only the auxiliary $V$ depending on $\pi$.
\item The matrix $u_Q$ intertwines the action of $\mathrm{GL}_n$ inside $\mathrm{GL}_{n+1}$ with the Iwahori-level-matching at conductor $\mathfrak q$; without it the integral would vanish identically when $c(\pi)\geq 1$.
\item An alternative construction uses the Godement--Jacquet approach: define $V$ as a translate of $V^\circ_\pi$ by an element of the Bernstein center of $\mathrm{GL}_n$.
\end{enumerate}

\section*{References}
\begin{itemize}\setlength\itemsep{0.2em}
\item H.~Jacquet, I.~Piatetski-Shapiro, J.~Shalika, \emph{Rankin-Selberg convolutions}, Amer.\ J.\ Math.\ 105 (1983), 367--464.
\item N.~Matringe, \emph{Essential Whittaker functions for $\mathrm{GL}(n)$}, Doc.\ Math.\ 18 (2013), 1191--1214.
\item N.~Matringe, \emph{Essential vectors and the $L$-function for $\mathrm{GL}_n\times\mathrm{GL}_m$}, IMRN (2013).
\item J.~Cogdell, I.~Piatetski-Shapiro, \emph{Converse theorems for $\mathrm{GL}_n$}, PSPM 61 (1997).
\end{itemize}

\end{document}
